\section{Work Plan}
\label{sec:plan}

\subsection{Completed}
\begin{itemize}
  \item Literature review of the reference study~\cite{severson2019} and its
  follow-up~\cite{attia2020}.
  \item Conversion of the 2018-02-20 batch from MATLAB to Python.
  \item Study of the dataset's structure, units and charging protocols.
  \item Cleaning (removal of cells without a label) and exploratory
  analysis.
  \item First computation of the $\dq$ feature and its correlation with
  cycle life.
  \item Problem statement finalised (Section~\ref{sec:problem}).
  \item \textbf{P1:} all four batches converted and cleaned (rules A1--A4,
  B1--B10); the 124 paper cells match the paper's Supplementary Table~9
  (Section~\ref{sec:pipeline}).
  \item \textbf{P2:} the 20 candidate features, the baselines and the
  paper's three models implemented; the paper's Table~1 reproduced
  (Section~\ref{sec:reproduction}).
  \item \textbf{P3:} 19 physics-informed candidates screened on the
  training cells; feature set F4 fixed at two features
  (Section~\ref{sec:physics}).
  \item \textbf{P4:} PLS, Gaussian process, SVR and additive splines
  compared with the elastic net; prediction error for 20--150 early cycles
  (Section~\ref{sec:models}).
  \item \textbf{P5:} frozen models applied to the unseen batch with five
  pre-registered hypotheses; SHAP interpretation (Section~\ref{sec:unseen}).
\end{itemize}

\subsection{Phases}
\begin{enumerate}
  \item \textbf{Data (P1).} Download the three batches of~\cite{matr_data}.
  Update the conversion notebooks to current \texttt{h5py} and apply the
  authors' published cleaning steps (removing faulty cells and joining
  cells that continued into the next batch). Save a lightweight cache
  holding the summaries and the \texttt{Qdlin}/\texttt{Tdlin} curves for
  cycles 1--100.
  \item \textbf{Baseline (P2).} Implement feature sets F1--F3, the
  mean and policy-mean baselines, and the elastic-net models; reproduce the
  errors reported in~\cite{severson2019} (RQ1).
  \item \textbf{Physics features (P3), done.} Screen physics-based
  candidates on the training cells; F4 = margin/rate and overhang
  rise (Section~\ref{sec:physics}).
  \item \textbf{Model comparison (P4), done.} PLS, Gaussian process, SVR and
  additive splines against the elastic net; error against the number of
  early cycles (Section~\ref{sec:models}).
  \item \textbf{Generalisation and interpretation (P5), done.} Unseen batch
  (RQ5) and SHAP analysis (Section~\ref{sec:unseen}).
  \item \textbf{Write-up (P6).} Final report and presentation.
\end{enumerate}

\subsection{Timeline}
\begin{table}[h]
\centering
\caption{Planned timeline (weeks from start of phase P1).}
\label{tab:timeline}
\begin{tabular}{@{}lll@{}}
\toprule
Weeks & Phase & Milestone \\
\midrule
1--2  & P1, P2 & Data cache ready; baseline reproduced \\
3--4  & P3     & Feature set F4 implemented \\
5--7  & P4     & Full results table; error-vs-$N$ plot \\
8     & P5     & Unseen-batch results; importance plots \\
9--10 & P6     & Final report and presentation \\
\bottomrule
\end{tabular}
\end{table}

\noindent\todo{Map weeks to calendar dates agreed with the guide.}

\subsection{Risks and mitigations}
\begin{itemize}
  \item \textbf{Large raw files (several GB per batch).} Convert each batch
  once and work from the lightweight cache.
  \item \textbf{Overfitting with about 124 samples.} Use regularisation,
  nested cross-validation and fixed test sets, and compare against simple
  baselines.
  \item \textbf{Physics features give no improvement.} Report this as a
  finding about which mechanisms do and do not predict lifetime.
\end{itemize}
